\section{闭合证书的复合、平稳投影与时间尺度边界}
\label{app:composition-time}

本附录分别处理两种累积：多个表示层依次粗粒化构成\emph{空间复合}，同一
近似因子反复演化构成\emph{时间迭代}。前者给出多层表示的误差复合律，后者
从给定的一步闭合证书推导固定时间与增长时间窗的充分条件。

\subsection{两层与多层空间复合}

设
\[
 Q:(X,K)\longrightarrow(Y,L),
 \qquad
 R:(Y,L)\longrightarrow(Z,M)
\]
是 Markov 通道，其中记号记录其源、靶动力学；交换缺陷由下式度量。

\begin{theorem}[近似动力学因子的复合]
\label{thm:app-factor-composition}
有
\begin{equation}
 \partial_{K,M}(QR)
 \le
 \alpha(R)\partial_{K,L}(Q)+\partial_{L,M}(R).
 \label{eq:app-two-layer-composition}
\end{equation}
特别地，若两个缺陷分别不超过 $\eps$ 与 $\delta$，则复合缺陷不超过
$\alpha(R)\eps+\delta\le\eps+\delta$。
\end{theorem}

\begin{proof}
由三角不等式和引理~\ref{lem:app-tv-contraction}，
\begin{align*}
 d_\infty(KQR,QRM)
 &\le d_\infty((KQ)R,(QL)R)
      +d_\infty(Q(LR),Q(RM))\\
 &\le \alpha(R)d_\infty(KQ,QL)+d_\infty(LR,RM).
\end{align*}
\end{proof}

称
\[
 (Q;\eps_Q,\rho_Q):(X,K)\longrightarrow(Y,L)
\]
为一个\emph{带收缩标签的近似因子证书}，若
\begin{equation}
 \partial_{K,L}(Q)\le\eps_Q,
 \qquad
 \alpha(Q)\le\rho_Q,
 \qquad
 \eps_Q\ge0,
 \quad \rho_Q\in[0,1].
 \label{eq:app-contraction-labelled-certificate}
\end{equation}
这里 $\rho_Q$ 控制表示通道 $Q$ 的 Dobrushin 系数；宏观动力学 $L$ 的混合率
另由 $\alpha(L)$ 度量。非恒定确定性通道满足 $\rho_Q=1$ 的最坏情形界。

若 $(Q;\eps_Q,\rho_Q)$ 与 $(R;\eps_R,\rho_R)$ 类型相接，则
定理~\ref{thm:app-factor-composition}和 Dobrushin 系数的次乘性说明，复合
$QR$ 可按
\begin{equation}
 (\eps_Q,\rho_Q)\star(\eps_R,\rho_R)
 :=(\rho_R\eps_Q+\eps_R,\rho_Q\rho_R)
 \label{eq:app-certificate-composition}
\end{equation}
复合。结合律来自普通乘法的结合律与
\[
 \rho_S(\rho_R\eps_Q+\eps_R)+\eps_S
 =(\rho_S\rho_R)\eps_Q+\rho_S\eps_R+\eps_S.
\]
保留 $(\eps,\rho)$ 标签可得到上述闭合复合律。若仅记录误差，复合仍满足加法
上界，而固定正误差阈值的未标记通道集合一般不对复合封闭。

\begin{corollary}[多层空间误差]
\label{cor:app-multilayer-composition}
对
\[
 Q_i:(X_{i-1},K_{i-1})\to(X_i,K_i),
 \qquad
 \eps_i:=\partial_{K_{i-1},K_i}(Q_i),
\]
令 $Q_{1:n}:=Q_1\cdots Q_n$。则
\begin{align}
 d_\infty(K_0Q_{1:n},Q_{1:n}K_n)
 &\le\sum_{i=1}^n
 \alpha(Q_{i+1}\cdots Q_n)\eps_i
 \label{eq:app-multilayer-exact-alpha}\\
 &\le\sum_{i=1}^n\eps_i
       \prod_{j=i+1}^n\alpha(Q_j)
 \le\sum_{i=1}^n\eps_i,
 \label{eq:app-multilayer-product-alpha}
\end{align}
其中空后缀按恒等通道处理，其误差权重取 $1$。
\end{corollary}

\begin{proof}
在有限维有符号核空间中作望远镜展开：
\[
 K_0Q_{1:n}-Q_{1:n}K_n
 =\sum_{i=1}^n Q_1\cdots Q_{i-1}
   (K_{i-1}Q_i-Q_iK_i)Q_{i+1}\cdots Q_n.
\]
前复合不增加行全变差范数，后复合至多乘以后缀的 Dobrushin 系数；三角
不等式给出式~\eqref{eq:app-multilayer-exact-alpha}。再用
$\alpha(Q_{i+1}\cdots Q_n)\le\prod_{j=i+1}^n\alpha(Q_j)$ 得到其余各式。
\end{proof}

\subsection{终端宏观边缘的多步误差}

\begin{theorem}[同一因子的 $n$ 步终端误差]
\label{thm:app-multistep-marginal}
若 $\partial_{K,L}(Q)\le\eps$，则对每个整数 $n\ge1$，
\begin{equation}
 d_\infty(K^nQ,QL^n)
 \le \eps\sum_{j=0}^{n-1}\alpha(L^j).
 \label{eq:app-multistep-general}
\end{equation}
若 $\alpha(L)\le\rho<1$，则
\begin{equation}
 d_\infty(K^nQ,QL^n)
 \le\eps\frac{1-\rho^n}{1-\rho}
 \le\frac{\eps}{1-\rho};
 \label{eq:app-multistep-contracting}
\end{equation}
不加收缩假设时仍有
\begin{equation}
 d_\infty(K^nQ,QL^n)\le n\eps.
 \label{eq:app-multistep-linear}
\end{equation}
这些界也控制任意初始分布 $\mu$ 的终端边缘
$\|\mu K^nQ-\mu QL^n\|_{\TV}$。
\end{theorem}

\begin{proof}
将行全变差范数延拓到行和为零的有符号核，使用恒等式
\[
 K^nQ-QL^n
 =\sum_{i=0}^{n-1}K^i(KQ-QL)L^{n-1-i}.
\]
每项的前复合不扩张，后复合至多乘以
$\alpha(L^{n-1-i})$；求和并重编号得到
式~\eqref{eq:app-multistep-general}。由
$\alpha(L^j)\le\alpha(L)^j$ 得到几何级数界；只用
$\alpha(L^j)\le1$ 则得到线性界。最后应用
式~\eqref{eq:app-mixture-bound}。
\end{proof}

式~\eqref{eq:app-multistep-contracting} 控制时间 $n$ 的\emph{终端宏观
分布}；整段路径的误差由定理~\ref{thm:app-path-coupling}中的顺序耦合另行
控制。

\subsection{近似平稳状态的投影}

\begin{proposition}[近似平稳状态投影]
\label{prop:stationary-projection}
设 $Q:X\rightsquigarrow Y$、$K:X\rightsquigarrow X$、
$L:Y\rightsquigarrow Y$，且
\[
d_\infty(KQ,QL)\le\eps,
\qquad
\|\mu K-\mu\|_{\TV}\le\beta.
\]
则 $\nu=\mu Q$ 满足
\[
\|\nu L-\nu\|_{\TV}\le\eps+\beta.
\]
若 $\alpha(L)\le\rho<1$，则 $L$ 有唯一平稳分布 $\pi_L$，且
\[
\|\mu Q-\pi_L\|_{\TV}
\le\frac{\eps+\beta}{1-\rho}.
\]
\end{proposition}

\begin{proof}
由混合凸性和推前非扩张，
\[
\|\mu QL-\mu Q\|_{\TV}
\le\|\mu QL-\mu KQ\|_{\TV}
+\|\mu KQ-\mu Q\|_{\TV}
\le\eps+\beta.
\]
若 $\alpha(L)\le\rho<1$，则 $\lambda\mapsto\lambda L$ 在有限概率单纯形上
严格压缩，因而存在唯一不动点。再由
\[
\|\nu-\pi_L\|_{\TV}
\le\|\nu-\nu L\|_{\TV}
+\|\nu L-\pi_LL\|_{\TV}
\le\eps+\beta+\rho\|\nu-\pi_L\|_{\TV}
\]
移项即得。这里的长期稳定性由近似平稳性与宏观收缩共同保证，群作用分离结构
在该命题中不承担条件作用。
\end{proof}

\subsection{确定性粗粒化的路径误差}

\begin{theorem}[宏观路径耦合]
\label{thm:app-path-coupling}
设 $q:X\twoheadrightarrow Y$，且
$\partial_{K,L}(Q_q)\le\eps\le1$。令 $(X_t)_{t\ge0}$ 是任意初始分布
$\mu$ 下的 $K$-链，并令 $Z_t=q(X_t)$。令 $(Y_t)_{t\ge0}$ 是初始分布
$\mu Q_q$、转移核 $L$ 的 Markov 链。则
\begin{equation}
 \left\|
 \operatorname{Law}(Z_0,\ldots,Z_n)
 -\operatorname{Law}(Y_0,\ldots,Y_n)
 \right\|_{\TV}
 \le1-(1-\eps)^n
 \le n\eps.
 \label{eq:app-path-bound}
\end{equation}
\end{theorem}

\begin{proof}
单步缺陷意味着对每个 $x$，
\[
 \left\|
 \operatorname{Law}(q(X_{t+1})\mid X_t=x)
 -L(q(x),\cdot)
 \right\|_{\TV}\le\eps.
\]
给定任一正概率宏观历史 $Z_{0:t}=z_{0:t}$，$X_t$ 的条件分布支撑在
$q^{-1}(z_t)$ 上。因此 $Z_{t+1}$ 的条件分布是上述纤维内分布的凸组合，
与 $L(z_t,\cdot)$ 仍相距至多 $\eps$。

从共同初态开始递归构造耦合：只要此前两条宏观路径相同，就对上述两个下一步
条件分布使用最大耦合，使下一状态继续相同的条件概率至少为 $1-\eps$；一旦
分离，按保持各自正确边缘的任意方式继续。经过 $n$ 次转移仍完全相同的概率
至少为 $(1-\eps)^n$。耦合不等式给出第一界，Bernoulli 不等式给出第二界。
\end{proof}

路径定理采用确定性 $Q$。对随机通道，方程 $KQ=QL$ 约束每一时刻的通道输出
边缘，过程级联合采样还需额外规定。例如取
$K=\Id_X$、$L=\Id_Y$，并令 $Q$ 的某行是非退化分布，则交换式精确成立。
若每个时刻独立地从 $Q(X_t,\cdot)$ 重采样，所得输出路径会变化，$L$-链则
保持不变。随机粗粒化的路径结论需要额外的联合耦合或 link consistency
条件。

\subsection{增长时间尺度的充分条件}

考虑随规模 $N$ 变化的证书
\[
 \partial_{K_N,L_N}(Q_N)\le\eps_N,
 \qquad \eps_N\longrightarrow0,
\]
以及观察时长 $T_N$。由前述定理可以严格区分三种情形：
\begin{enumerate}
 \item 对每个固定 $T$，$\eps_N\to0$ 足以使终端误差和确定性粗粒化的路径
  误差趋零。
 \item 对确定性粗粒化的完整路径，一个不依赖额外结构的直接充分条件是
 \begin{equation}
   T_N\eps_N\longrightarrow0.
   \label{eq:app-growing-horizon-linear}
 \end{equation}
 对终端边缘，则有更精确的充分条件
 \begin{equation}
   \eps_N\sum_{j=0}^{T_N-1}\alpha(L_N^j)\longrightarrow0.
   \label{eq:app-growing-horizon-alpha}
 \end{equation}
 \item 若存在与 $N$ 无关的 $\rho<1$ 使 $\alpha(L_N)\le\rho$，则
 $\eps_N/(1-\rho)\to0$ 控制所有 $T_N$ 的终端宏观边缘；整段路径继续由
 式~\eqref{eq:app-path-bound} 控制，因为路径距离记录早期偏差。
\end{enumerate}

增长时长的控制需要误差率与时间窗之间的联合条件。一个最小示例是
$X=Y=\{0,1\}$、$Q=\Id$、$L=\Id$，而 $K_\eps$ 从 $0$ 以概率
$\eps$ 跳到吸收态 $1$、从 $1$ 保持不动。此时
\[
 d_\infty(K_\eps Q,QL)=\eps,
\]
从 $0$ 出发，在 $n$ 步后与 $L$-链终端分布的全变差距离为
\[
 1-(1-\eps)^n.
\]
取 $n=\lfloor1/\eps\rfloor$，该量不趋于零。这个例子固定 $Q$ 与 $L$，刻画
给定近似因子证书的传播边界；联合优化表示与 $L$ 属于另一问题，在本例中取
$L=K_\eps$ 可得到零缺陷。

正文的系统族主定理给出一步结论，并直接推出固定有限时间结论。增长时间窗的
终端边缘由式~\eqref{eq:app-growing-horizon-alpha}、统一混合或其他明确稳定性
条件控制；完整路径由式~\eqref{eq:app-growing-horizon-linear}或更强的路径级
耦合条件控制。宏观收缩改善终端边缘界，路径距离则保留已经发生的早期偏差。
$\Gamma_N\to0$ 与这些时间尺度条件共同使用时，才能得到相应增长时间结论。
